\subsection*{Construcción de la evidencia y validación}
Se emplearon 197 parcelas y un conjunto supervisado de 138 rendimientos. BASIC reúne observaciones multitemporales Sentinel y Landsat; PRO incorpora observaciones Planet. Se agregaron indicadores del ciclo de 2025, antecedentes de vegetación, lluvia, topografía y geología con comprobaciones de identificadores y trazabilidad. CHIRPS, SoilGrids y WaPOR se tratan como fuentes indirectas con escala propia \parencite{Funk2015CHIRPS,Poggio2021SoilGrids,FAOWaPOR,Dhonthi2024}. La validación consistió en 138 exclusiones sucesivas de una parcela, utilizando las otras 137 etiquetas para cada ajuste. Las 59 etiquetas ocultas permanecieron inaccesibles. Una transformación sin dependencia del rendimiento puede consultar las 197 covariables; toda operación ajustada con rendimientos utiliza únicamente etiquetas disponibles en su entrenamiento \parencite{Stone1974,Chapelle2006}.

\subsection*{Estrategias de imputación}
Se consideraron mediana y media aritmética calculadas dentro de cada conjunto de entrenamiento, así como reconstrucción espacial, temporal y espaciotemporal de predictores faltantes. Las estrategias temporales respetaron fechas, sensores y calidad de observación; las espaciales exploraron parcelas próximas y contexto territorial. No se confundió completar una señal espectral con estimar el rendimiento. La versión final emplea media porque obtuvo un RMSE ligeramente inferior al de la mediana en una comparación de 138 exclusiones. No se atribuye superioridad frente a interpolaciones alternativas sin su tabla final auditada ni se sostiene que una diferencia mínima sea significativa.

\subsection*{Predicción climática y experto local}
El componente C--T3 aplica Elastic Net sobre cuatro entradas: lluvia total del ciclo, fracción correspondiente a meses tempranos, concentración mensual de precipitación y NDVI Sentinel--2 en mayo de 2025 \parencite{ZouHastie2005}. El experto local usa siete atributos distintos que resumen nivel y variabilidad histórica de NDVI, variación fenológica observada por Planet, estadísticos de NDWI y pendiente. Para cada parcela selecciona los 24 vecinos geográficamente más cercanos con rendimiento conocido, ajusta una regresión Ridge ponderada inversamente por distancia y conserva penalización 30 \parencite{HoerlKennard1970}. Un ajuste residual interno Elastic Net, basado en predicciones honestas del experto, complementa la estimación local.

\subsection*{Ponderación geológica y residuales}
Las diez proporciones de composición geológica, transformadas mediante raíz cuadrada y centrado, alimentan un K-means de cuatro u ocho grupos; la alternativa se elige dentro del entrenamiento. Cada grupo recibe un peso regularizado que mezcla las predicciones global y local. Las coordenadas sirven para elegir vecinos del experto y no se confunden con las variables de agrupamiento geológico \parencite{MacQueen1967}. Se ensayó DBSCAN como alternativa de agrupamiento por densidad, con manejo explícito de ruido y asignación de parcelas nuevas, sin cambiar los demás componentes \parencite{Ester1996DBSCAN}. El corrector R3 utiliza tres indicadores de estructura histórica del NDVI, persistencia de heterogeneidad y pendiente mínima; es regularizado, se desactiva si faltan indicadores y está acotado. En conjunto, los 26 atributos se distribuyen en cuatro globales, siete locales, dos coordenadas, diez geológicos y tres residuales.

\subsection*{Corrección final y elección interna}
Se examinó la cercanía entre parcelas en 26 atributos escalados robustamente y, alternativamente, en ocho componentes principales. La regla seleccionada considera los tres análogos etiquetados más próximos, ponderados inversamente por su distancia en atributos, y corrige parcialmente la predicción base con un factor de mezcla de 0.25. Una compuerta reduce el cambio cuando la cercanía o la concordancia de rendimientos vecinos indica soporte débil; el ajuste absoluto se limita a 0.5 toneladas por hectárea. La elección entre cinco recetas finalistas y la ausencia de corrección se analizó también mediante cuatro particiones internas en cada exclusión externa. Aunque este anidamiento evita que la etiqueta del caso retenido intervenga directamente en la elección interna, los candidatos habían sido preseleccionados tras explorar esas mismas 138 etiquetas.
